%!TEX root =  tfg.tex

\begin{quotation}[Empresario]{Steve Jobs (1955--2011)}
El diseño no es solo lo que ves, sino cómo funciona.
\end{quotation}

\begin{abstract}
En este capítulo reflejaremos todos y cada uno de los costes que se han de reflejar a la hora de desarrollar el proyecto así como la relación de funcionalidades que se desarrollan con estos costes.
\end{abstract}

\section{Características a desarrollar}

\begin{enumerate}
\item Funcionalidad 1: Búsqueda de todas las tareas. Ver Tabla \ref{tab:valorAportado1Iteracion1}.
\item Funcionalidad 2: Creación de tareas. Ver Tabla \ref{tab:valor2Iteracion1}.
\item Funcionalidad 3: Eliminación de tareas. Ver Tabla \ref{tab:valor3Iteracion1}.
\item Funcionalidad 4: Asignación de estados a las tareas. Ver Tabla \ref{tab:valor4Iteracion1}.
\item Funcionalidad 5: Búsqueda de tareas acabadas e inacabadas. Ver Tabla \ref{tab:valor5Iteracion1}.
\item Funcionalidad 6: Encontrar compañeros de un equipo. Ver Tabla \ref{tab:valor6Iteracion1}.
\end{enumerate}
%----------------------------------------------------------------------------------------------------------
\begin{table*}[htb]
	\centering
	\begin{coolTable}{p{4cm}p{\textwidth-4.5cm}}{2}
{Análisis de valor aportado 0001}
\textbf{Propuesta}&Se pretende desarrollar la funcionalidad de listar todas las tareas creadas por los usuarios\\
\midrule
\textbf{Valor}&El valor aportado es poder visualizar las tareas creadas.\\
\textbf{Coste}&Se requiere la participación de ambos miembros del equipo desarrollador para llevar a cabo la implementación.\\
\textbf{Opciones}&Esta funcionalidad tiene un valor prioritario en la aplicación.\\
\textbf{Riesgos}&La no correcta visualización de los diferentes campos que tiene una tarea.\\
\textbf{Deuda técnica}&Se ha elegido la solución más factible para el desarrollo de la misma.\\
	\end{coolTable}
	\caption{Análisis de valor aportado 0001\label{tab:valorAportado1Iteracion1}}
\end{table*}
%----------------------------------------------------------------------------------------------------------
\begin{table*}[htb]
	\centering
	\begin{coolTable}{p{4cm}p{\textwidth-4.5cm}}{2}
{Análisis de valor aportado 0002}
\textbf{Propuesta}&Se pretende desarrollar la funcionalidad de crear tareas por los usuarios.\\
\midrule
\textbf{Valor}&El valor aportado es poder crear las tareas por los usuarios.\\
\textbf{Coste}&Se requiere la participación de ambos miembros del equipo desarrollador para llevar a cabo la implementación.\\
\textbf{Opciones}&Esta funcionalidad tiene un valor prioritario en la aplicación.\\
\textbf{Riesgos}&La posibilidad de no poder crear la tarea de forma satisfactoria.\\
\textbf{Deuda técnica}&Se ha elegido la solución más factible para el desarrollo de la misma.\\
	\end{coolTable}
	\caption{Análisis de valor aportado 0002\label{tab:valor2Iteracion1}}
\end{table*}
%----------------------------------------------------------------------------------------------------------
\begin{table*}[htb]
	\centering
	\begin{coolTable}{p{4cm}p{\textwidth-4.5cm}}{2}
{Análisis de valor aportado 0003}
\textbf{Propuesta}&Se pretende desarrollar la funcionalidad de eliminar tareas por los usuarios.\\
\midrule
\textbf{Valor}&El valor aportado es poder eliminar las tareas por los usuarios.\\
\textbf{Coste}&Se requiere la participación de ambos miembros del equipo desarrollador para llevar a cabo la implementación.\\
\textbf{Opciones}&Esta funcionalidad tiene un valor prioritario en la aplicación.\\
\textbf{Riesgos}&La posibilidad de no poder eliminar las tarea de forma satisfactoria.\\
\textbf{Deuda técnica}&Se ha elegido la solución más factible para el desarrollo de la misma.\\
	\end{coolTable}
	\caption{Análisis de valor aportado 0003\label{tab:valor3Iteracion1}}
\end{table*}

%----------------------------------------------------------------------------------------------------------
\begin{table*}[htb]
	\centering
	\begin{coolTable}{p{4cm}p{\textwidth-4.5cm}}{2}
{Análisis de valor aportado 0004}
\textbf{Propuesta}&Se pretende desarrollar la funcionalidad de asignar estado a las tareas por los usuarios.\\
\midrule
\textbf{Valor}&El valor aportado es poder asignar a las tareas el estado en el que se encuentren por los usuarios.\\
\textbf{Coste}&Se requiere la participación de ambos miembros del equipo desarrollador para llevar a cabo la implementación.\\
\textbf{Opciones}&Esta funcionalidad tiene un valor alto en la aplicación.\\
\textbf{Riesgos}&La posibilidad de no poder asignar el estado a una tarea.\\
\textbf{Deuda técnica}&Se ha elegido la solución más factible para el desarrollo de la misma.\\
	\end{coolTable}
	\caption{Análisis de valor aportado 0004\label{tab:valor4Iteracion1}}
\end{table*}

%----------------------------------------------------------------------------------------------------------
\begin{table*}[htb]
	\centering
	\begin{coolTable}{p{4cm}p{\textwidth-4.5cm}}{2}
{Análisis de valor aportado 0005}
\textbf{Propuesta}&Se pretende desarrollar la funcionalidad de listar todas las tareas acabadas e inacabadas por los usuarios\\
\midrule
\textbf{Valor}&El valor aportado es poder visualizar las tareas acabadas e inacabadas por los usuarios.\\
\textbf{Coste}&Se requiere la participación de ambos miembros del equipo desarrollador para llevar a cabo la implementación.\\
\textbf{Opciones}&Esta funcionalidad tiene un valor medio en la aplicación.\\
\textbf{Riesgos}&La no correcta visualización de los estados en las tareas creadas por los usuarios.\\
\textbf{Deuda técnica}&Se ha elegido la solución más factible para el desarrollo de la misma.\\
	\end{coolTable}
	\caption{Análisis de valor aportado 0005\label{tab:valor5Iteracion1}}
\end{table*}

%----------------------------------------------------------------------------------------------------------
\begin{table*}[htb]
	\centering
	\begin{coolTable}{p{4cm}p{\textwidth-4.5cm}}{2}
{Análisis de valor aportado 0006}
\textbf{Propuesta}&Se pretende desarrollar la funcionalidad de encontrar los compañeros de un equipo.\\
\midrule
\textbf{Valor}&El valor aportado es poder visualizar los compañeros de un equipo de trabajo.\\
\textbf{Coste}&Se requiere la participación de ambos miembros del equipo desarrollador para llevar a cabo la implementación.\\
\textbf{Opciones}&Esta funcionalidad tiene un valor medio en la aplicación.\\
\textbf{Riesgos}&La no correcta visualización de los compañeros de un equipo.\\
\textbf{Deuda técnica}&Se ha elegido la solución más factible para el desarrollo de la misma.\\
	\end{coolTable}
	\caption{Análisis de valor aportado 0006\label{tab:valor6Iteracion1}}
\end{table*}

\clearpage

\section{Diseño}
Aquí una discusión de cómo va a afectar todo al diseño

Debe insertarse un diagrama UML de diseño con los cambios y hacer referencia en el texto así Fig. \ref{fig:diseno01}.

\begin{figure}[htbp]
\begin{center}

\includegraphics[width=\textwidth]{figures/UML-iteracion-1.jpg}
\caption{Diagrama UML de diseño para la iteración 1}
\label{fig:diseno01}
\end{center}
\end{figure}

\clearpage
Un memorando técnico por cada decisión de diseño.

\begin{table*}[htb]
	\centering
	\begin{coolTable}{p{4cm}p{\textwidth-4.5cm}}{2}
{Memorando técnico 0001}
\textbf{Asunto}&Diseño del modelo Tarea.\\
\textbf{Resumen}&Se ha diseñado un modelo de Tarea que contenga los elementos necesarios para nuestra definición de esta; título, descripción, fechas de inicio y fin, duración, priorización si está terminada o no, la estimación de la finalización de la tarea, además contaremos con el creador de la tarea, el usuario asignado a la tarea y el equipo de la tarea si esta tuviera.\\
\midrule
\textbf{Factores causantes}&La relación entre modelos, ya que la clase Tarea cuenta con una asociación tanto con el modelo Equipo como con el modelo Usuario, así como los datos que requerimos que estos tengan asignados para el proyecto. \\
\textbf{Solución}&El modelo resultante asigna un tipo de dato para cada uno de los campos así como las relaciones que no pueden ser nulas en ningún momento por la concepción del modelo de Tarea como pueden ser el nombre o el creador.\\
\textbf{Motivación}&Un modelo que contenga los datos necesarios para un correcto procesamiento de las tareas.\\
\textbf{Cuestiones abiertas}& La correcta relación entre tablas y completitud de los atributos.\\
\textbf{Alternativas}&Se valoraron otras alternativas de relaciones, pero la dispuesta es la que se consideró más óptima tanto a niveles de implementación como de seguridad.\\
	\end{coolTable}
	\caption{Memorando técnico 0001}
\end{table*}

%-----------------------------------------------------------------------------------------------------------
\begin{table*}[htb]
	\centering
	\begin{coolTable}{p{4cm}p{\textwidth-4.5cm}}{2}
{Memorando técnico 0002}
\textbf{Asunto}&Diseño del modelo Equipo.\\
\textbf{Resumen}&Se ha diseñado un modelo de Equipo que contenga los elementos necesarios para nuestra definición de esta; nombre, además contaremos con el creador del equipo, lista de asociados, así como de la lista de managers.\\
\midrule
\textbf{Factores causantes}&La relación entre modelos, ya que la clase Equipo cuenta con una asociación con el modelo Usuario, así como los datos que requerimos que estos tengan asignados para el proyecto.\\
\textbf{Solución}&El modelo resultante asigna un tipo de dato para cada uno de los campos.\\
\textbf{Motivación}&Un modelo que contenga los datos necesarios para un correcto procesamiento de los equipos.\\
\textbf{Cuestiones abiertas}& La correcta relación entre tablas y completitud de los atributos.\\
\textbf{Alternativas}&Dado a la tecnología utilizada la manera de llevar a cabo esta relación entre tablas es la más óptima.\\
	\end{coolTable}
	\caption{Memorando técnico 0002}
\end{table*}

%-------------------------------------------------------------------------------------------------------

\begin{table*}[htb]
	\centering
	\begin{coolTable}{p{4cm}p{\textwidth-4.5cm}}{2}
{Memorando técnico 0003}
\textbf{Asunto}&Diseño del modelo Usuario.\\
\textbf{Resumen}&Se ha diseñado un modelo de Usuario que contenga los elementos necesarios para nuestra definición de este; nombre, apellidos, correo electrónico, la hora de inicio de su jornada laboral, hora de fin de la jornada laboral, nombre de usuario, contraseña y una lista de equipos.\\
\midrule
\textbf{Factores causantes}&La relación entre modelos, ya que la clase Usuario cuenta con dos asociaciones con el modelo Tarea y con el modelo Equipo, así como los datos que requerimos que estos tengan asignados para el proyecto.\\
\textbf{Solución}&El modelo resultante asigna un tipo de dato para cada uno de los campos.\\
\textbf{Motivación}&Un modelo que contenga los datos necesarios para un correcto procesamiento de las Usuario.\\
\textbf{Cuestiones abiertas}&La correcta relación entre tablas y completitud de los atributos.\\
\textbf{Alternativas}&Se valoraron otras alternativas de relaciones, pero la dispuesta es la que se consideró más óptima tanto a niveles de implementación como de seguridad.\\
	\end{coolTable}
	\caption{Memorando técnico 0003}
\end{table*}

%----------------------------------------------------------------------------------------------------------
\begin{table*}[htb]
	\centering
	\begin{coolTable}{p{4cm}p{\textwidth-4.5cm}}{2}
{Memorando técnico 0004}
\textbf{Asunto}&Diseño del modelo Invitación.\\
\textbf{Resumen}&Se ha diseñado un modelo de Invitación que contenga los elementos necesarios para nuestra definición de esta; mensaje, estado, correo electrónico, además contaremos con el usuario invitado, así como el equipo al que pertenece la invitación.\\
\midrule
\textbf{Factores causantes}&La relación entre modelos, ya que la clase Invitación cuenta con una asociación con el modelo Usuario y con una asociación con el modelo Equipo, así como los datos que requerimos que estos tengan asignados para el proyecto.\\
\textbf{Solución}&El modelo resultante asigna un tipo de dato para cada uno de los campos.\\
\textbf{Motivación}&Un modelo que contenga los datos necesarios para un correcto procesamiento de las invitaciones.\\
\textbf{Cuestiones abiertas}& La correcta relación entre tablas y completitud de los atributos.\\
\textbf{Alternativas}&Dado a la tecnología utilizada la manera de llevar a cabo esta relación entre tablas es la más óptima.\\
	\end{coolTable}
	\caption{Memorando técnico 0004}
\end{table*}

%------------------------------------------------------------------------------------------------------

\clearpage

\section{Implementación}

En esta sección se va a realizar un memorando técnico por cada decisión de implementación y refactorización que afecte al diseño.

\begin{table*}[htb]
	\centering
	\begin{coolTable}{p{4cm}p{\textwidth-4.5cm}}{2}
{Memorando técnico 0001}
\textbf{Asunto}&La creación de la clase Tarea.\\
\textbf{Resumen}&Se han definido los siguientes atributos para la clase Tarea que son los necesarios para mostrar los datos que queremos.\\
\midrule
\textbf{Factores causantes}&La relación entre tablas, ya que la clase Tarea cuenta con una asociación tanto con la clase Equipo como con la clase Usuario. \\
\textbf{Solución}&La solución llevada a cabo tiene la utilización de las anotaciones que proporciona JPA, para gestionar las relaciones entre dos tablas. En este caso para la relación entre Tarea y Equipo será \textbf{@ManyToOne} y para la relación entre Tarea y Usuario será \textbf{@OneToOne}.\\
\textbf{Motivación}&Ya conocíamos esta forma de relacionar las tablas.\\
\textbf{Cuestiones abiertas}& La correcta relación entre tablas.\\
\textbf{Alternativas}&Dado a la tecnología utilizada la manera de llevar a cabo esta relación entre tablas es la más óptima.\\
	\end{coolTable}
	\caption{Memorando técnico 0001}
\end{table*}

%-----------------------------------------------------------------------------------------------------------
\begin{table*}[htb]
	\centering
	\begin{coolTable}{p{4cm}p{\textwidth-4.5cm}}{2}
{Memorando técnico 0002}
\textbf{Asunto}&La creación de la clase Equipo.\\
\textbf{Resumen}&Se han definido los siguientes atributos para la clase Equipo que son los necesarios para mostrar los datos que queremos.\\
\midrule
\textbf{Factores causantes}&La relación entre tablas, ya que la clase Equipo cuenta con una asociación con la clase Usuario. \\
\textbf{Solución}&La solución llevada a cabo tiene la utilización de las anotaciones que proporciona JPA, para gestionar las relaciones entre dos tablas. En este caso la relación entre Equipo y Usuario serán \textbf{@ManyToMany}.\\
\textbf{Motivación}&Ya conocíamos esta forma de relacionar las tablas.\\
\textbf{Cuestiones abiertas}& La correcta relación entre tablas.\\
\textbf{Alternativas}&Dado a la tecnología utilizada la manera de llevar a cabo esta relación entre tablas es la más óptima.\\
	\end{coolTable}
	\caption{Memorando técnico 0002}
\end{table*}

%-------------------------------------------------------------------------------------------------------

\begin{table*}[htb]
	\centering
	\begin{coolTable}{p{4cm}p{\textwidth-4.5cm}}{2}
{Memorando técnico 0003}
\textbf{Asunto}&La creación de la clase Usuario.\\
\textbf{Resumen}&Se han definido los siguientes atributos para la clase Usuario que son los necesarios para mostrar los datos que queremos.\\
\midrule
\textbf{Factores causantes}&La relación entre tablas, ya que la clase Usuario cuenta con una asociación con la clase Equipo. \\
\textbf{Solución}&La solución llevada a cabo tiene la utilización de las anotaciones que proporciona JPA, para gestionar las relaciones entre dos tablas. En este caso para la relación entre Usuario y Equipo será \textbf{@OneToMany}.\\
\textbf{Motivación}&Ya conocíamos esta forma de relacionar las tablas.\\
\textbf{Cuestiones abiertas}& La correcta relación entre tablas.\\
\textbf{Alternativas}&Dado a la tecnología utilizada la manera de llevar a cabo esta relación entre tablas es la más óptima.\\
	\end{coolTable}
	\caption{Memorando técnico 0003}
\end{table*}

%----------------------------------------------------------------------------------------------------------
\begin{table*}[htb]
	\centering
	\begin{coolTable}{p{4cm}p{\textwidth-4.5cm}}{2}
{Memorando técnico 0004}
\textbf{Asunto}&La creación de la clase Invitación.\\
\textbf{Resumen}&Se han definido los siguientes atributos para la clase Invitación que son los necesarios para mostrar los datos que queremos.\\
\midrule
\textbf{Factores causantes}&La relación entre tablas, ya que la clase Invitación cuenta con una asociación tanto con la clase Equipo como con la clase Usuario. \\
\textbf{Solución}&La solución llevada a cabo tiene la utilización de las anotaciones que proporciona JPA, para gestionar las relaciones entre dos tablas. En este caso para la relación entre Invitación y Equipo será \textbf{@OneToOne} y para la relación entre Invitación y Usuario será \textbf{@OneToOne}.\\
\textbf{Motivación}&Ya conocíamos esta forma de relacionar las tablas.\\
\textbf{Cuestiones abiertas}& La correcta relación entre tablas.\\
\textbf{Alternativas}&Dado a la tecnología utilizada la manera de llevar a cabo esta relación entre tablas es la más óptima.\\
	\end{coolTable}
	\caption{Memorando técnico 0004}
\end{table*}

%-----------------------------------------------------------------------------------------------------

\clearpage

\section{Pruebas}

De la misma manera que con la iteración anterior no se han integrado ninguna prueba automatizada sobre el código, estas se desarrollarán en su totalidad en la sexta iteración. Por lo que por ahora aunque no se han desarrollado pruebas automatizadas, el software producido ha sido testeado por los desarrolladores de manera manual.

\section{Despliegue}

Se han realizado constantes despliegues en la iteración 1, al menos se han realizado despliegues cada dos días desde el comienzo de la iteración salvo el último fin de semana de la misma en la que los despliegues fueron diarios para las últimas pruebas de la obtención de datos.